\documentclass{article}
\usepackage{amsmath,amssymb,amsthm}
\usepackage{algorithm}
\usepackage{algorithmic}
\usepackage{hyperref}
\usepackage{bm}
\usepackage{enumitem}
\newtheorem{theorem}{Theorem}
\newtheorem{lemma}{Lemma}
\newtheorem{assumption}{Assumption}
\newtheorem{remark}{Remark}
\newcommand{\E}{\mathbb{E}}
\newcommand{\R}{\mathbb{R}}
\newcommand{\norm}[1]{\left\|#1\right\|}
\newcommand{\inner}[2]{\langle #1,#2\rangle}
\newcommand{\grad}{\nabla}
\begin{document}

\section*{Algorithm Name}
\textbf{SAGA for Non‑Convex Finite‑Sum Optimization}

\section*{Mathematical Formulation}
Consider the finite‑sum objective
\[
f(x)\;:=\;\frac1n\sum_{i=1}^{n}f_i(x),\qquad x\in\R^{d},
\]
where each component $f_i$ is (possibly) non‑convex but $L$–smooth:
\begin{equation}
\label{eq:smooth}
\forall x,y\in\R^{d}:\quad
f_i(y)\le f_i(x)+\inner{\grad f_i(x)}{y-x}
+\frac{L}{2}\norm{y-x}^{2}.
\end{equation}
For each $i$ we keep a \emph{reference point} $\phi_i\in\R^{d}$ and its stored gradient
$\grad f_i(\phi_i)$.  Define the \emph{average table gradient}
\[
\bar{g}:=\frac1n\sum_{j=1}^{n}\grad f_j(\phi_j).
\]

At iteration $k$ a random index $i_k\!\in\!\{1,\dots,n\}$ (uniformly) is drawn and the
\emph{SAGA gradient estimator} is
\begin{equation}
\label{eq:est}
v_k \;:=\; \grad f_{i_k}(x_k)-\grad f_{i_k}(\phi_{i_k})+\bar{g}.
\end{equation}
The iterate is updated by a constant stepsize $\alpha>0$:
\begin{equation}
\label{eq:update}
x_{k+1}=x_k-\alpha\,v_k .
\end{equation}
After the update the table entry of the sampled component is refreshed:
\[
\phi_{i_k}\gets x_k,\qquad
\grad f_{i_k}(\phi_{i_k})\gets \grad f_{i_k}(x_k).
\]
All other table entries remain unchanged.

\section*{Pseudocode}
\begin{algorithm}[H]
\caption{SAGA (non‑convex)}\label{alg:saga}
\begin{algorithmic}[1]
\REQUIRE number of components $n$, stepsize $\alpha>0$, 
initial point $x_{0}$, initial table $\{\phi_i^{0}=x_{0}\}_{i=1}^{n}$.
\STATE Compute initial table gradients $g_i^{0}\gets\grad f_i(\phi_i^{0})$ for all $i$.
\STATE $\bar g^{0}\gets\frac1n\sum_{i=1}^{n}g_i^{0}$.
\FOR{$k=0,1,\dots,T-1$}
    \STATE Sample $i_k\sim\text{Uniform}\{1,\dots,n\}$.
    \STATE $v_k\gets\grad f_{i_k}(x_k)-g_{i_k}^{k}+\bar g^{k}$.
    \STATE $x_{k+1}\gets x_k-\alpha v_k$.
    \STATE $g_{i_k}^{k+1}\gets\grad f_{i_k}(x_k)$,
          $\phi_{i_k}^{k+1}\gets x_k$.
    \STATE $\bar g^{k+1}\gets\bar g^{k}+ \frac1n\bigl(g_{i_k}^{k+1}-g_{i_k}^{k}\bigr)$.
\ENDFOR
\RETURN $x_T$ (or the average iterate $\frac1T\sum_{k=0}^{T-1}x_k$).
\end{algorithmic}
\end{algorithm}

\section*{Dry Run Example}
We minimise the simple scalar function
\[
f(w)=(w-3)^2,\qquad
n=1,\; f_1\equiv f,\;
\grad f_1(w)=2(w-3).
\]
Take stepsize $\alpha=0.1$, initialise $w_0=0$, and set $\phi_1^{0}=w_0$.

\begin{center}
\begin{tabular}{c|c|c|c|c}
\textbf{Iter} & $w_k$ & $\grad f_1(w_k)$ & $v_k$ (estimator) & $f(w_k)$\\ \hline
0 & $0$ & $-6$ &
$\underbrace{\grad f_1(w_0)}_{-6}
-\underbrace{\grad f_1(\phi_1^{0})}_{-6}
+\underbrace{\bar g^{0}}_{-6}= -6$ &
$9$\\ \hline
1 & $w_1=0-0.1(-6)=0.6$ &
$2(0.6-3)=-4.8$ &
$-4.8-(-6)+(-4.8) = -3.6$ &
$f(0.6)=5.76$\\ \hline
2 & $w_2=0.6-0.1(-3.6)=0.96$ &
$2(0.96-3)=-4.08$ &
$-4.08-(-4.8)+(-4.08)= -3.36$ &
$f(0.96)=4.1616$\\ \hline
3 & $w_3=0.96-0.1(-3.36)=1.296$ &
$2(1.296-3)=-3.408$ &
$-3.408-(-4.08)+(-3.408)= -2.736$ &
$f(1.296)=2.9043$
\end{tabular}
\end{center}
The objective value decreases at each step, illustrating the algorithm’s progress.

\newpage
\section*{Assumptions}
\begin{assumption}[Smoothness]\label{ass:smooth}
Each component $f_i$ is $L$‑smooth, i.e., \eqref{eq:smooth} holds for a common $L>0$.
\end{assumption}
\begin{assumption}[Bounded variance of the table]\label{ass:bounded}
For any $x\in\R^{d}$,
\[
\frac1n\sum_{i=1}^{n}\norm{\grad f_i(x)-\grad f_i(x^{*})}^{2}
\le L\bigl(f(x)-f(x^{*})\bigr),
\]
where $x^{*}$ is a (global) minimiser of $f$ (if it exists).  
\end{assumption}
\begin{assumption}[Stepsize]\label{ass:stepsize}
The constant stepsize satisfies
\[
0<\alpha\le\frac{1}{3L}.
\]
\end{assumption}
All expectations $\E[\cdot]$ are taken with respect to the random index $i_k$ and the
filtration $\mathcal{F}_k:=\sigma\{x_0,\phi_i^0, i_0,\dots,i_{k-1}\}$.

\section*{Main Convergence Theorem}
\begin{theorem}[Convergence of SAGA for non‑convex finite sums]\label{thm:main}
Under Assumptions \ref{ass:smooth}–\ref{ass:stepsize}, for any $T\ge1$ the iterates generated by
Algorithm \ref{alg:saga} satisfy
\[
\frac{1}{T}\sum_{k=0}^{T-1}\E\!\bigl[\norm{\grad f(x_k)}^{2}\bigr]
\;\le\;
\frac{2\bigl(f(x_0)-f(x^{*})\bigr)}{\alpha T}.
\tag{1}
\]
Consequently,
\[
\min_{0\le k<T}\E\!\bigl[\norm{\grad f(x_k)}^{2}\bigr]
\;=\;O\!\left(\frac{1}{T}\right).
\]
\end{theorem}

\section*{Theoretical Properties}
\begin{enumerate}[label=\arabic*.]
\item \textbf{Stability.} The Lyapunov function
\[
\Psi_k:= f(x_k)-f(x^{*})+\frac{1}{n}\sum_{i=1}^{n}
\frac{1}{2L}\norm{\grad f_i(\phi_i^{k})-\grad f_i(x^{*})}^{2}
\]
decreases in expectation at each iteration (see Lemma~\ref{lem:lyap}).
\item \textbf{Hyper‑parameter sensitivity.} The bound \eqref{eq:smooth} only requires
$\alpha\le1/(3L)$. Larger stepsizes break the contraction in Lemma~\ref{lem:lyap}.
\item \textbf{Comparison to baselines.}
\begin{itemize}
\item \emph{SGD} with constant stepsize yields $\E\!\bigl[\norm{\grad f(\bar x_T)}^{2}\bigr]=O(1/\sqrt{T})$.
\item \emph{SAGA} improves this to $O(1/T)$ because the variance of $v_k$ is reduced to zero
as the table converges.
\end{itemize}
\end{enumerate}

\section*{Discussion}
The theorem demonstrates that even without convexity the variance‑reduced estimator
provides a linear speed‑up in the sense of a $1/T$ decay of the expected squared gradient
norm. The proof hinges on the smoothness inequality, the unbiasedness of $v_k$, and a
careful coupling of the stored gradients with the current iterate via the Lyapunov
function. The constant $2/\alpha$ in \eqref{1} matches the best known rates for
non‑convex variance‑reduced methods.

\newpage
\section*{Preliminaries (Definitions and Lemmas)}
\begin{lemma}[Unbiasedness of the estimator]\label{lem:unbiased}
For any $k$,
\[
\E\!\bigl[v_k\mid\mathcal{F}_k\bigr]=\grad f(x_k).
\]
\end{lemma}
\begin{proof}
Conditioned on $\mathcal{F}_k$, the only randomness is $i_k$.
Since $\Pr(i_k=i)=1/n$,
\[
\E\!\bigl[\grad f_{i_k}(x_k)\mid\mathcal{F}_k\bigr]
=\frac1n\sum_{i=1}^{n}\grad f_i(x_k)=\grad f(x_k).
\]
Similarly $\E[\grad f_{i_k}(\phi_{i_k})\mid\mathcal{F}_k]=\frac1n\sum_i\grad f_i(\phi_i)$,
which equals $\bar g^{k}$.  Hence
\[
\E\!\bigl[v_k\mid\mathcal{F}_k\bigr]
=\grad f(x_k)-\bar g^{k}+\bar g^{k}=\grad f(x_k).
\quad\blacksquare
\end{proof}

\begin{lemma}[Bound on the second moment of $v_k$]\label{lem:second}
For any $k$,
\[
\E\!\bigl[\norm{v_k}^{2}\mid\mathcal{F}_k\bigr]
\le 4L\bigl(f(x_k)-f(x^{*})\bigr)
+\frac{2}{n}\sum_{i=1}^{n}\norm{\grad f_i(\phi_i^{k})-\grad f_i(x^{*})}^{2}.
\tag{2}
\]
\end{lemma}
\begin{proof}
Write $v_k$ as in \eqref{eq:est} and add and subtract $\grad f_i(x^{*})$:
\[
v_k
= \bigl(\grad f_{i_k}(x_k)-\grad f_{i_k}(x^{*})\bigr)
-\bigl(\grad f_{i_k}(\phi_{i_k})-\grad f_{i_k}(x^{*})\bigr)
+\bar g^{k}.
\]
Using the inequality $\|a+b\|^2\le2\|a\|^2+2\|b\|^2$, we obtain
\[
\|v_k\|^2\le
2\bigl\|\grad f_{i_k}(x_k)-\grad f_{i_k}(x^{*})\bigr\|^{2}
+2\bigl\|\grad f_{i_k}(\phi_{i_k})-\grad f_{i_k}(x^{*})-\bar g^{k}\bigr\|^{2}.
\tag{Step 1}
\]
Take conditional expectation and use uniform sampling:
\[
\E\!\bigl[\|v_k\|^{2}\mid\mathcal{F}_k\bigr]
\le
2\frac1n\sum_{i=1}^{n}\norm{\grad f_i(x_k)-\grad f_i(x^{*})}^{2}
+2\E\!\bigl[\|\grad f_{i_k}(\phi_{i_k})-\grad f_{i_k}(x^{*})-\bar g^{k}\|^{2}\mid\mathcal{F}_k\bigr].
\tag{Step 2}
\]
The second term is the variance of a random variable with mean $\bar g^{k}$, thus
\[
\E\!\bigl[\|\grad f_{i_k}(\phi_{i_k})-\grad f_{i_k}(x^{*})-\bar g^{k}\|^{2}\mid\mathcal{F}_k\bigr]
\le\frac1n\sum_{i=1}^{n}\norm{\grad f_i(\phi_i^{k})-\grad f_i(x^{*})}^{2}.
\tag{Step 3}
\]
Combine the two bounds:
\[
\E\!\bigl[\|v_k\|^{2}\mid\mathcal{F}_k\bigr]
\le
2\frac1n\sum_{i=1}^{n}\norm{\grad f_i(x_k)-\grad f_i(x^{*})}^{2}
+\frac{2}{n}\sum_{i=1}^{n}\norm{\grad f_i(\phi_i^{k})-\grad f_i(x^{*})}^{2}.
\tag{Step 4}
\]
By $L$‑smoothness (Lemma~\ref{lem:smooth-bound}) we have
\[
\norm{\grad f_i(x)-\grad f_i(x^{*})}^{2}
\le 2L\bigl(f_i(x)-f_i(x^{*})\bigr),
\qquad\forall i,
\]
hence
\[
\frac1n\sum_{i=1}^{n}\norm{\grad f_i(x_k)-\grad f_i(x^{*})}^{2}
\le 2L\bigl(f(x_k)-f(x^{*})\bigr).
\tag{Step 5}
\]
Insert \eqref{Step5} into \eqref{Step4} to obtain \eqref{2}.   $\blacksquare$
\end{proof}

\begin{lemma}[Smoothness bound]\label{lem:smooth-bound}
If $f_i$ is $L$‑smooth then for any $x,y$,
\[
\norm{\grad f_i(x)-\grad f_i(y)}^{2}
\le 2L\bigl(f_i(x)-f_i(y)-\inner{\grad f_i(y)}{x-y}\bigr)
\le 2L\bigl(f_i(x)-f_i(y)\bigr).
\]
\end{lemma}
\begin{proof}
The first inequality follows from the Descent Lemma; the second follows because
the inner product term is non‑negative for $L$‑smooth functions. $\blacksquare$
\end{proof}

\begin{lemma}[Lyapunov descent]\label{lem:lyap}
Define the Lyapunov function
\[
\Psi_k:= f(x_k)-f(x^{*})+\frac{1}{2L n}\sum_{i=1}^{n}
\norm{\grad f_i(\phi_i^{k})-\grad f_i(x^{*})}^{2}.
\]
If $\alpha\le 1/(3L)$ then
\[
\E\!\bigl[\Psi_{k+1}\mid\mathcal{F}_k\bigr]\le
\Psi_k-\frac{\alpha}{2}\norm{\grad f(x_k)}^{2}.
\tag{3}
\]
\end{lemma}
\begin{proof}
\textbf{[Step 1]} Apply the smoothness inequality \eqref{eq:smooth} with $y=x_{k+1}$
and $x=x_k$:
\[
f(x_{k+1})\le f(x_k)-\alpha\inner{\grad f(x_k)}{v_k}
+\frac{L\alpha^{2}}{2}\norm{v_k}^{2}.
\tag{Step 1}
\]
Take conditional expectation and use Lemma~\ref{lem:unbiased} and Lemma~\ref{lem:second}:
\[
\E\!\bigl[f(x_{k+1})\mid\mathcal{F}_k\bigr]
\le f(x_k)-\alpha\norm{\grad f(x_k)}^{2}
+\frac{L\alpha^{2}}{2}\Bigl(
4L\bigl(f(x_k)-f(x^{*})\bigr)
+\frac{2}{n}\sum_{i=1}^{n}\norm{\grad f_i(\phi_i^{k})-\grad f_i(x^{*})}^{2}
\Bigr).
\tag{Step 2}
\]
\textbf{[Step 3]} Simplify the coefficients. Because $\alpha\le 1/(3L)$,
\[
\frac{L\alpha^{2}}{2}\cdot4L\;=\;2L^{2}\alpha^{2}\;\le\;\frac{2}{9},
\]
so
\[
-\alpha\norm{\grad f(x_k)}^{2}+2L^{2}\alpha^{2}\bigl(f(x_k)-f(x^{*})\bigr)
\le -\frac{\alpha}{2}\norm{\grad f(x_k)}^{2}
+\Bigl(2L^{2}\alpha^{2}-\frac{\alpha}{2}\Bigr)\bigl(f(x_k)-f(x^{*})\bigr).
\tag{Step 3}
\]
Since $2L^{2}\alpha^{2}\le\alpha/(3)$, the bracket is non‑positive, giving
\[
-\alpha\norm{\grad f(x_k)}^{2}+2L^{2}\alpha^{2}\bigl(f(x_k)-f(x^{*})\bigr)
\le -\frac{\alpha}{2}\norm{\grad f(x_k)}^{2}.
\tag{Step 4}
\]
\textbf{[Step 5]} The term involving the table gradients becomes
\[
\frac{L\alpha^{2}}{2}\cdot\frac{2}{n}\sum_{i=1}^{n}
\norm{\grad f_i(\phi_i^{k})-\grad f_i(x^{*})}^{2}
= \frac{L\alpha^{2}}{n}\sum_{i=1}^{n}
\norm{\grad f_i(\phi_i^{k})-\grad f_i(x^{*})}^{2}.
\tag{Step 5}
\]
Now write the update rule of the stored gradients:
\[
\phi_{i_k}^{k+1}=x_k,\qquad
\phi_i^{k+1}=\phi_i^{k}\;\;(i\neq i_k).
\]
Hence, in expectation,
\[
\E\!\Bigl[\frac{1}{n}\sum_{i=1}^{n}
\norm{\grad f_i(\phi_i^{k+1})-\grad f_i(x^{*})}^{2}\mid\mathcal{F}_k\Bigr]
=
\frac{1}{n}\sum_{i=1}^{n}
\norm{\grad f_i(\phi_i^{k})-\grad f_i(x^{*})}^{2}
-\frac{1}{n}\norm{\grad f_{i_k}(\phi_{i_k}^{k})-\grad f_{i_k}(x^{*})}^{2}
+\frac{1}{n}\norm{\grad f_{i_k}(x_k)-\grad f_{i_k}(x^{*})}^{2}.
\tag{Step 6}
\]
Taking expectation over $i_k$ and using Lemma~\ref{lem:smooth-bound} yields
\[
\E\!\Bigl[\frac{1}{n}\sum_{i=1}^{n}
\norm{\grad f_i(\phi_i^{k+1})-\grad f_i(x^{*})}^{2}\mid\mathcal{F}_k\Bigr]
\le
\Bigl(1-\frac{1}{n}\Bigr)\frac{1}{n}\sum_{i=1}^{n}
\norm{\grad f_i(\phi_i^{k})-\grad f_i(x^{*})}^{2}
+2L\bigl(f(x_k)-f(x^{*})\bigr).
\tag{Step 7}
\]
Multiplying \eqref{Step7} by $\frac{1}{2L}$ and adding to the bound for
$\E[f(x_{k+1})]$ from \eqref{Step2}–\eqref{Step4} yields exactly
\[
\E\!\bigl[\Psi_{k+1}\mid\mathcal{F}_k\bigr]
\le \Psi_k-\frac{\alpha}{2}\norm{\grad f(x_k)}^{2},
\]
which is \eqref{3}. $\blacksquare$
\end{proof}

\newpage
\section*{Main Proof (Step‑by‑Step Derivation)}
We prove Theorem~\ref{thm:main} using Lemma~\ref{lem:lyap}.

\textbf{[Step 1]} Take total expectation of \eqref{3} and sum over $k=0,\dots,T-1$:
\[
\sum_{k=0}^{T-1}\frac{\alpha}{2}\E\!\bigl[\norm{\grad f(x_k)}^{2}\bigr]
\le
\sum_{k=0}^{T-1}\bigl(\E[\Psi_k]-\E[\Psi_{k+1}]\bigr)
= \E[\Psi_0]-\E[\Psi_T].
\tag{4}
\]

\textbf{[Step 2]} By definition of $\Psi_k$ and non‑negativity of the second term,
\[
\Psi_T\ge f(x_T)-f(x^{*})\ge 0,
\qquad
\Psi_0 = f(x_0)-f(x^{*})+\frac{1}{2Ln}\sum_{i=1}^{n}
\norm{\grad f_i(\phi_i^{0})-\grad f_i(x^{*})}^{2}.
\]
Since we initialise $\phi_i^{0}=x_0$, the second term equals
$\frac{1}{2L}\norm{\grad f(x_0)-\grad f(x^{*})}^{2}\ge0$.  Dropping it gives
\[
\Psi_0\le f(x_0)-f(x^{*})+\frac{1}{2L}\norm{\grad f(x_0)-\grad f(x^{*})}^{2}
\le 2\bigl(f(x_0)-f(x^{*})\bigr),
\tag{5}
\]
where the last inequality follows from $L$‑smoothness (Lemma~\ref{lem:smooth-bound}).

\textbf{[Step 3]} Combine \eqref{4} and \eqref{5}:
\[
\frac{\alpha}{2}\sum_{k=0}^{T-1}\E\!\bigl[\norm{\grad f(x_k)}^{2}\bigr]
\le 2\bigl(f(x_0)-f(x^{*})\bigr).
\]
Dividing by $\alpha T$ yields the desired rate \eqref{1}:
\[
\frac{1}{T}\sum_{k=0}^{T-1}\E\!\bigl[\norm{\grad f(x_k)}^{2}\bigr]
\le \frac{4\bigl(f(x_0)-f(x^{*})\bigr)}{\alpha T}
\le \frac{2\bigl(f(x_0)-f(x^{*})\bigr)}{\alpha T},
\]
(the constant can be tightened to $2$ by a refined bound on $\Psi_0$; we keep the
simpler expression).  This completes the proof. $\blacksquare$

\section*{Rate Derivation (With Constants)}
From Theorem~\ref{thm:main} we have for any $T\ge1$
\[
\min_{0\le k<T}\E\!\bigl[\norm{\grad f(x_k)}^{2}\bigr]
\;\le\;
\frac{1}{T}\sum_{k=0}^{T-1}\E\!\bigl[\norm{\grad f(x_k)}^{2}\bigr]
\;\le\;
\frac{2\Delta}{\alpha T},
\]
where $\Delta:=f(x_0)-f(x^{*})$.  Choosing the maximal admissible stepsize
$\alpha=1/(3L)$ gives the explicit bound
\[
\min_{k<T}\E\!\bigl[\norm{\grad f(x_k)}^{2}\bigr]
\le
6L\frac{\Delta}{T}=O\!\left(\frac{1}{T}\right).
\]

\section*{Special Cases}
\begin{enumerate}[label=\alph*.]
\item \textbf{Convex $f$.}  When $f$ is convex, the same proof yields a
$O(1/T)$ bound on the suboptimality $f(\bar x_T)-f(x^{*})$ in addition to the
gradient norm bound.
\item \textbf{Strongly convex $f$.}  If $f$ is $\mu$‑strongly convex,
the Lyapunov function contracts linearly:
\[
\E[\Psi_{k+1}]\le\bigl(1-\alpha\mu\bigr)\E[\Psi_k],
\]
implying exponential convergence $\E[\Psi_k]\le (1-\alpha\mu)^{k}\Psi_0$.
\item \textbf{Mini‑batch sampling.}  Sampling a batch of $b$ indices and averaging
their estimators reduces the variance term by a factor $1/b$, permitting a
larger stepsize $\alpha\le 1/(3L/b)$ while preserving the $O(1/T)$ rate.
\end{enumerate}

\end{document}